\section{Detailed Design of the Processing Pipeline}

This section specifies each of the four pipeline stages in turn: the models employed,
the parameters governing them, and the output each stage produces for the next. Every
numeric value given is the value in the implementation.

% =============================================================================
\subsection{Stage 1 --- Chunking}

\subsubsection*{Purpose and inputs}

The stage receives a decoded video file and produces an ordered list of non-overlapping
time spans covering the recording. Audio is decoded once with PyAV into a mono waveform,
from which the duration is derived; the video stream is read separately by the visual
detectors.

\subsubsection*{Boundary detectors}

Four detectors run once per recording. Each returns a list of
(time, strength) events, where the strength lies in $[0,1]$.

\begin{table}[htbp]
\centering
\small
\begin{tabularx}{\textwidth}{@{}L{1.9cm}L{3.5cm}XL{2.5cm}@{}}
\toprule
\textbf{Signal} & \textbf{Model / algorithm} & \textbf{Event derivation} & \textbf{Strength} \\
\midrule
Speaker change & \texttt{pyannote/speaker-\allowbreak diarization-\allowbreak community-1} (gated, requires a Hugging Face token) &
The waveform is diarized into speaker turns; a boundary is emitted at the midpoint
between two consecutive turns whose speaker label differs. &
Constant $1.0$. A detected speaker change is a binary, high-confidence event. \\
Silence & Silero VAD v6 &
Speech timestamps are returned for the whole track; a boundary is emitted at the
midpoint of every gap between speech segments longer than $0.05$\,s, including the
trailing gap. &
$\min(1,\, g/2.0)$ for a gap of $g$ seconds, saturating at two seconds. \\
Hard cut & PySceneDetect \texttt{ContentDetector}; purely algorithmic frame differencing, no model weights &
A boundary is emitted at the start of every scene after the first. &
Constant $1.0$. The detector applies its own threshold, so any report is confident. \\
Semantic drift & OpenCLIP ViT-B/32, OpenAI weights, frames sampled at 1\,fps &
Consecutive sampled frames are embedded and compared; a boundary is emitted at each
midpoint with raw strength $1-\cos(e_i, e_{i+1})$. &
Rescaled by the clip's own largest jump, so the strongest semantic change is placed on
equal footing with a speaker change or a cut. \\
\bottomrule
\end{tabularx}
\end{table}

The rescaling applied to the semantic signal is necessary because CLIP similarity drifts
continuously and rarely approaches zero even at a genuine scene change; without it the
signal would be permanently diluted relative to the sparse, unit-strength detectors.

Frames for the semantic detector are read by walking the container forward with
\texttt{grab()}/\texttt{retrieve()} rather than seeking to each timestamp. Per-frame
seeking forces the H.264 decoder to restart from a keyframe and was measured at
approximately $22\times$ the cost per frame.

\subsubsection*{Fusion}

Weights are first renormalised over the signals that produced at least one event, so
that a preset expresses \emph{which signals matter} rather than acting as a granularity
dial (Section~4, E-1). Each event then contributes a Gaussian bump to a score curve
evaluated on a grid of resolution $0.1$\,s over the recording:

\[
S(t) \;=\; \sum_{s \in S} \; \sum_{(t_i,\, a_i) \in E_s}
w_s \, a_i \, \exp\!\left(-\frac{(t - t_i)^2}{2\sigma^2}\right),
\qquad \sigma = 1.0\ \mathrm{s}
\]

Here $S$ is the set of signals, $E_s$ the events emitted by signal $s$, $w_s$ its
renormalised weight, $a_i$ the strength of event $i$, and $\sigma$ the width of the
Gaussian kernel.

Boundaries agreed upon by several detectors therefore reinforce one another. Peaks are
selected greedily in descending order of $S(t)$, subject to a score floor of $0.12$ and
a minimum separation of $2.0$\,s between accepted boundaries, which prevents runs of
adjacent micro-chunks.

\begin{table}[htbp]
\centering
\small
\begin{tabularx}{\textwidth}{@{}L{2.6cm}cccc X@{}}
\toprule
\textbf{Preset} & \textbf{speaker} & \textbf{silence} & \textbf{cut} & \textbf{semantic} & \textbf{Intended footage} \\
\midrule
\texttt{audio}       & 0.35 & 0.35 & 0.15 & 0.15 & Dialogue-led recordings \\
\texttt{video}       & 0.15 & 0.15 & 0.35 & 0.35 & Edited or visually varied footage \\
\texttt{audio\_video} & 0.25 & 0.25 & 0.25 & 0.25 & Default; all four signals equally \\
\bottomrule
\end{tabularx}
\end{table}

\subsubsection*{Modes and post-processing}

Three mutually exclusive modes are exposed. In \texttt{preset} mode a named weighting is
used; in \texttt{weights} mode the caller supplies its own, validated against the four
signal names and required to be non-zero in at least one; in \texttt{interval} mode the
recording is cut into fixed spans of $N$ seconds and no detector runs at all.

For the two weighted modes, a span of length $d$ longer than the maximum duration
$D_{\max}$ is divided into $\lceil d / D_{\max} \rceil$ equal parts --- equal division
rather than truncation, so that no negligible trailing fragment is produced --- and a
span shorter than the minimum duration is merged forward into its successor, with a
short final span folded back into the preceding one. The API defaults are $5$\,s and
$20$\,s.
Duration bounds are deliberately \emph{not} applied in \texttt{interval} mode, since the
caller requested exact spans.

\subsubsection*{Output}

An ordered list of (start, end) spans, together with a chunk
configuration key that identifies the scheme that produced them --- for example
\texttt{audio\_video:5-20} or \texttt{interval:10}. Custom weightings are hashed into
the key so that two different weightings sharing the same duration bounds cannot
collide. This key is carried on every downstream vector, which is what permits two
chunkings of one recording to coexist.

% =============================================================================
\subsection{Stage 2 --- Analysis}

\subsubsection*{Purpose and inputs}

Each selected analyzer receives the chunk spans and a video context that provides
frames and audio on demand. Every analyzer owns its own frame sampling, its own prompt
and its own output schema, and is registered by identifier; nothing in ingestion,
storage or the API names a specific analyzer.

\subsubsection*{Common sampling strategy}

The frame-selecting analyzers share a three-part strategy. Candidate frames are drawn
at a low fixed rate; near-duplicates are removed by comparing OpenCLIP embeddings
against the last frame retained; and the survivors are thinned to a per-analyzer
maximum. The scene analyzer derives its similarity threshold adaptively, as the
40th percentile of consecutive-frame similarity within the chunk, capped at
$0.98$. The cap is load-bearing: on a static chunk the percentile lands on the median,
measured at $0.9955$, and distinct frames are then mined out of compression noise. The
detector-gated analyzers instead use a fixed similarity threshold together with a
layout-change test, because a person or object entering or leaving the frame is a
change the whole-image comparison does not register.

All frames selected for a chunk are sent in a single request, so that the model sees
the chunk as a continuous sequence and can describe movement across it rather than
judging each frame in isolation. Sampling is serial, because the CLIP model is shared
on one GPU; the API calls are I/O-bound and run across six worker threads.

\subsubsection*{Analyzers, models and outputs}

\begin{table}[htbp]
\centering
\small
\begin{tabularx}{\textwidth}{@{}L{2.2cm}L{3.5cm}X@{}}
\toprule
\textbf{Analyzer} & \textbf{Models and gating} & \textbf{Structured output} \\
\midrule
\texttt{default\_video} &
Candidates at 2\,fps; OpenCLIP ViT-B/32 deduplication (adaptive threshold, cap $0.98$);
4--8 frames retained, thinned against evenly spaced points in \emph{time}; VLM
\texttt{gpt-5.4-mini} at 768\,px. &
Schema \texttt{scene\_description}: \texttt{description} (the prose users search
against), \texttt{setting}, \texttt{people[]}, \texttt{objects[]}, \texttt{actions[]},
\texttt{tags[]}. The timestamps of the frames used are recorded alongside for
traceability. \\
\texttt{people} &
YOLO11n gate at 1\,fps, confidence $0.35$; person boxes below $900$\,px$^2$ discarded;
identifiers assigned by overlap tracking across \emph{all} candidate frames, then read
off for the frames actually sent; up to 6 frames annotated with numbered yellow boxes;
VLM at 1280\,px, \texttt{detail=high}. &
Schema \texttt{people\_reading}: per person \texttt{box\_id}, \texttt{appearance},
\texttt{clothing}, \texttt{role}, \texttt{action}, plus a \texttt{summary}. Each person
is post-processed with \texttt{locations} --- the timestamp and box in every frame they
appear in --- and a \texttt{people\_count} is recorded. \\
\texttt{object\_\allowbreak detection} &
Shares the same YOLO11n pass rather than repeating it; frames annotated with cyan
\emph{unlabelled} boxes; up to 6 frames; VLM at 1280\,px, \texttt{detail=high}. &
Schema \texttt{object\_reading}: \texttt{detections[]} of \texttt{object},
\texttt{description} (colour, material, size, condition) and \texttt{context} (what it
is being used for, and by whom), plus a \texttt{summary}. Plain object names are
extracted into a filterable list; the detector's own labels are retained only for
debugging. \\
\texttt{ocr} &
EasyOCR \emph{detection} stage only --- its recognition stage is skipped, as it
fragments lines and reads low-contrast overlays poorly; frames annotated with red
boxes; similarity threshold $0.985$ plus a text-layout change test; VLM at 1600\,px,
\texttt{detail=high}. &
Schema \texttt{ocr\_reading}: \texttt{texts[]} of \texttt{text} (verbatim, with case and
punctuation) and \texttt{context} (what the text is and where it appears), plus a
\texttt{summary}. \\
\texttt{transcript} &
faster-whisper (\texttt{tiny}); language detected once over the whole track and pinned
for every chunk. &
One plain-text transcription per chunk. \\
\texttt{diarization} &
pyannote over the whole track once, memoised; faster-whisper segments per chunk;
speaker assigned by largest temporal overlap between a Whisper segment and a pyannote
turn. &
\texttt{turns[]} of \texttt{speaker}, \texttt{start}, \texttt{end}, \texttt{text} with
consecutive same-speaker segments merged, plus a rendered \texttt{text}, a
speaker-free \texttt{plain} form, and the \texttt{speakers} present. \\
\bottomrule
\end{tabularx}
\end{table}

The vision--language model is invoked with OpenAI strict structured output, so every
required property is present and no additional properties are returned; downstream
stages therefore never parse prose.

Two design rules apply across the table. First, whole-track passes are performed once
and memoised: a per-chunk diarization pass would label the same person
\texttt{SPEAKER\_00} in one chunk and \texttt{SPEAKER\_01} in the next, and per-segment
language detection produced fluent German from an English clip. Second,
\texttt{transcript} and \texttt{diarization} are mutually exclusive and the API enforces
this, because their embedded vectors were measured cosine-identical at $1.0000$ ---
diarization is a strict superset.

\subsubsection*{Gating and cost}

Where an analyzer's gate finds nothing --- no person above the area threshold, no object
detection, no text region --- the chunk yields no output and no API call is made at all.
A frame that survives is downscaled to the analyzer's width and JPEG-encoded before
transmission, since vision tokens are priced by resolution.

% =============================================================================
\subsection{Stage 3 --- Indexing and Aggregation}

\subsubsection*{Embedding}

Every analyzer implements \texttt{render\_fields}, which flattens its output into a map
of named vector space to text. Five spaces are defined --- \texttt{combined},
\texttt{description}, \texttt{people}, \texttt{actions}, \texttt{objects} --- of which
\texttt{combined} must always be present. A space is \emph{omitted} rather than emitted
empty, so a chunk in which the model saw nobody is simply not a candidate in a
people-scoped search.

Text is embedded locally with \texttt{BAAI/bge-small-en-v1.5} through
sentence-transformers, producing normalised vectors compared by cosine distance. The
model family was trained with an instruction prefix on the query side only, so queries
are prefixed with \emph{``Represent this sentence for searching relevant passages:''}
and documents are not; applying the prefix to both, or to neither, measurably degrades
recall. All fields of all chunks are embedded in one batched GPU pass and then
scattered back to their (chunk, field) slots.

\subsubsection*{Vector store}

One Qdrant point is written per chunk per analyzer, carrying all five named vectors.
Because the point is the unit, a limit of ten returns ten chunks, the payload is stored
once, and a filter is applied once. The point identifier is a deterministic UUID derived
from
\texttt{sha1(video\_id\,|\,analyzer\_id\,|\,start\,|\,end\,|\,chunk\_config)}, keyed on
the time span rather than a positional index --- chunk 12 of one chunking run is a
different moment from chunk 12 of another, so a positional identifier would silently
rebind a vector to the wrong timeframe. Re-ingesting the same chunk therefore upserts
rather than duplicating.

The payload carries everything required to cite and render a moment without joining back
to disk: the video identifier and playback URL, the analyzer, the chunk configuration
and identifier, the time span, the combined text and description, the setting, the
people, objects, actions and tags, the speakers and their turns, the read texts, the
object detections, the full per-person records, and the person count. Twelve payload
fields are declared as indexes, and twelve named filters are defined in a single filter
specification --- video, chunk and analyzer identifiers, chunk configuration, objects,
tags, speakers, people, minimum and maximum person count, and time bounds. Filters are
validated against that specification and an unknown name raises an error with the
nearest suggestion, because a silently ignored filter returns plausible but wrong
results.

\subsubsection*{Aggregators}

Twelve aggregators run over the saved analyzer records, in an order derived from their
declared dependencies. A dependency naming an analyzer is a requirement on the
recording, and an aggregator whose analyzer is absent is dropped rather than run against
nothing. Five of the twelve incur API calls; the remainder are free, which is what makes
re-running them on demand practical.

\begin{table}[htbp]
\centering
\small
\begin{tabularx}{\textwidth}{@{}L{2.3cm}L{2.6cm}L{2.9cm}X@{}}
\toprule
\textbf{Aggregator} & \textbf{Depends on} & \textbf{Model} & \textbf{Output} \\
\midrule
\texttt{stats} & --- & None; arithmetic &
Duration, chunk count, person count series with minimum, maximum, mean, busiest and
quietest spans; object frequencies; word count, spoken seconds and speech ratio. \\
\texttt{novelty} & --- & None; reuses stored vectors &
Cosine distance of each chunk's \texttt{combined} vector from the normalised centroid of
the recording, ranked; outliers are those beyond two standard deviations, so
``unusual'' scales with how varied the recording itself is. \\
\texttt{speaker\_stats} & \texttt{diarization} & None; arithmetic &
Per speaker: seconds, turns, words, share of speech, longest turn and chunks. Plus
handover count, dominant speaker and a turn timeline. \\
\texttt{summary} & --- & LLM &
A hierarchy built by repeated halving. Depth is derived from chunk count --- enough
halvings for a leaf to cover at most five chunks, bounded to between three and six tiers
--- so that leaf granularity does not vary with recording length. Leaves are summarised
from chunk text and every tier above merges its children, making the whole-recording
summary a reduction rather than one enormous prompt. Emits \texttt{summary},
\texttt{key\_points[]}, \texttt{topics[]}, the full tier hierarchy, and the finest tier
as \texttt{sections} for retrieval. \\
\texttt{chapters} & --- & LLM &
Runs of consecutive chunks grouped into named sections with a title and one-sentence
summary. Chapters must be consecutive, non-overlapping and exhaustive; any chapter
referencing a chunk identifier that does not exist is discarded. \\
\texttt{events} & --- & LLM &
Discrete, chunk-anchored occurrences: \texttt{event}, \texttt{chunk\_id},
\texttt{actor}, \texttt{category}, resolved to real time spans and counted by category.
Ongoing background activity is explicitly excluded. \\
\texttt{ner} & --- & GLiNER \texttt{urchade/\allowbreak gliner\_small-v2.1}, zero-shot, threshold $0.5$ &
Named entities over descriptions, on-screen text and speech across ten configurable
labels. Pronouns and bare determiners are excluded by a stop list, as they are
grammatically entities but identify nobody and would dominate a mention count. Emits
entities with mention counts, scores and chunk identifiers, grouped by label. \\
\texttt{sentiment} & \texttt{diarization} (falls back to \texttt{transcript}) & \texttt{cardiffnlp/\allowbreak twitter-roberta-\allowbreak base-sentiment-latest} &
A per-turn timeline, per-speaker distribution with dominant label and negative share,
and an overall distribution. Run on speech only, and reported as sentiment of language
observed rather than emotion inferred. \\
\texttt{entities} & \texttt{people} & BGE embeddings and constrained clustering, then LLM narratives &
Person observations linked into individuals. See below. \\
\texttt{object\_entities} & \texttt{object\_detection} & The same clusterer, plus lexical filters &
Object sightings linked into individual objects. Static fixtures are counted but not
tracked, and person-like detections are handed to the people aggregator. \\
\texttt{entity\_timelines} & \texttt{entities} & None &
Per person: first and last sighting, total span, \emph{observed} presence summed over
sighting spans rather than first-to-last, appearance count and spans. \\
\texttt{cooccurrence} & \texttt{entities} & None &
Pairs of people sharing a chunk with their counts, people per chunk, and the objects
present. \\
\bottomrule
\end{tabularx}
\end{table}

\subsubsection*{Entity linking in detail}

Identity is decided by embedding and clustering, not by asking a model to match people.
Measured on this footage, provably different people --- co-visible within one chunk, so
they cannot be the same --- reached $0.941$ similarity, while known-same pairs sat
between $0.889$ and $0.915$. Those bands overlap, so no threshold alone is sufficient
and the constraints carry the weight. The procedure is:

\begin{enumerate}[leftmargin=1.6em]
  \item \textbf{Signature construction.} The signature is the \texttt{clothing} field
        alone. Blending in \texttt{appearance} dragged same-person scores down to
        $0.69$--$0.87$, because that field drifts between chunks and occasionally
        carries meta-commentary that the embedding treats as content.
  \item \textbf{Within-chunk merge.} Observations in the same chunk above a similarity
        of $0.95$ --- above the measured co-visible ceiling --- are merged first. The
        analyzer's own tracker loses a person who moves between sampled frames, and
        without this step the cannot-link constraint would permanently forbid reuniting
        them.
  \item \textbf{Distinctness gate.} Each observation's mean similarity to all others is
        computed. Anything at or above $0.80$ describes half the recording and
        identifies nobody, so its owner is left unlinked rather than merged into
        whichever cluster it drifted closest to.
  \item \textbf{Constrained agglomeration.} Pairs above $0.88$, drawn from different
        chunks and both distinctive, are merged greedily in descending order of
        similarity. A merge is rejected if the two clusters share any chunk, since that
        would place one person in two places at once.
  \item \textbf{Narration.} Only then is an LLM used, and only for people seen more
        than once, capped at twelve narratives per recording. It produces a one-line
        recognition description, an ordered account of what the person did, and an
        apparent role.
\end{enumerate}

Object linking reuses the same routine on a source where identity is inherently weaker
--- two shopping carts genuinely look alike where two people rarely wear the same
clothes --- so the distinctness gate does most of the work.

% =============================================================================
\subsection{Stage 4 --- Retrieval and Question Answering}

\subsubsection*{Search}

A search embeds the query with the query-side prefix, selects one of the five named
vector spaces, applies the validated filter set, and returns ranked chunks with their
time spans and payloads. Two parameters deserve note.

\emph{Score threshold.} Cosine similarity always ranks something, so a query for content
a recording does not contain still returns its nearest neighbours. A score threshold is
what makes ``no match'' expressible. On this data, content actually present scored
$\geq 0.665$ and absent content $\leq 0.524$, so a threshold of approximately
$0.55$--$0.60$ separates them; the value is specific to the embedding model and is
retuned if the model changes.

\emph{Detail level.} Three levels are offered --- \texttt{minimal}, \texttt{standard}
and \texttt{full}. Five full hits were measured at approximately $19{,}000$ tokens
against approximately $440$ at \texttt{minimal}. The interface requests the whole
record; an agent, which pays per token, requests less.

\subsubsection*{Question routing}

A question is not answered by searching chunks and hoping. Each aggregate carries a
hint --- a short description, written in the vocabulary a question would use, of what
that aggregate is good for. The hints are embedded and compared against the question,
and the aggregates standing at least one standard deviation above the mean score for
that question are selected, capped at three. If nothing stands out, the single
best-matching aggregate is used, so a broad question is answered from the strongest
match rather than from nothing.

Relative selection is used because absolute thresholds do not survive contact with real
questions: the score distribution shifts per question, with means between $0.47$ and
$0.56$ measured here, so any single cutoff either admits everything or nothing. The
ranking, however, is reliable. Because routing is by embedded description rather than by
keyword, a new aggregator joins by adding one line of hint text, and unanticipated
phrasings degrade gracefully.

\begin{table}[htbp]
\centering
\small
\begin{tabularx}{\textwidth}{@{}L{2.9cm}X@{}}
\toprule
\textbf{Aggregate} & \textbf{Questions it is routed} \\
\midrule
\texttt{entities}        & What a particular person did; someone identified by clothing or appearance; following one individual; how long a person stayed. \\
\texttt{novelty}         & Which segment is unusual, anomalous or out of the ordinary; what stands out. \\
\texttt{stats}           & How many people were present; how busy or crowded; the busiest or quietest moment; counts and totals. \\
\texttt{events}          & What happened and when; arrivals, departures and transactions in time order. \\
\texttt{chapters}        & The structure of the footage; what sections it has; what the recording is about overall. \\
\texttt{speaker\_stats}  & Who talked most; how many speakers; speaking time, turns and interruptions. \\
\texttt{sentiment}       & The tone of what was said aloud; whether spoken language was positive, negative or neutral. \\
\texttt{ner}             & Named things mentioned: brands, organisations, places, names, dates and amounts. \\
\texttt{object\_entities}& A particular object followed through the footage; what an item was used for over time. \\
\texttt{cooccurrence}    & Which people appeared together; who was accompanied by whom. \\
\bottomrule
\end{tabularx}
\end{table}

\subsubsection*{Context assembly and synthesis}

The routed aggregates are rendered into labelled blocks, each carrying explicit
timecodes. The overall summary is always included, together with the four summary
sections closest to the question and, where entities were routed, the four people whose
descriptions best match it. Matching chunks from vector search are appended last. The
response records which aggregates were routed to and how much of each was used, so an
answer can be traced to its sources.

The synthesis prompt constrains the model to the supplied material, requires every claim
to carry a \texttt{[video\_id start-end]} citation drawn from the timecodes given, and
directs it to state plainly that the material does not support an answer rather than
guessing (NFR-5).

\subsubsection*{Output to the client}

The stage returns cited answers with timecodes, ranked moments with playable time ranges
and playback URLs, and --- from the aggregates --- chapters, the event timeline, entity
timelines, co-occurrence, statistics, novelty and sentiment. A moment is delivered as a
range over a recording rather than as a rendered file, so several moments drawn from one
recording share a single media load in the client.
